% Options for packages loaded elsewhere
% Options for packages loaded elsewhere
\PassOptionsToPackage{unicode}{hyperref}
\PassOptionsToPackage{hyphens}{url}
\PassOptionsToPackage{dvipsnames,svgnames,x11names}{xcolor}
%
\documentclass[
  british,
]{article}
\usepackage{xcolor}
\usepackage{amsmath,amssymb}
\setcounter{secnumdepth}{5}
\usepackage{iftex}
\ifPDFTeX
  \usepackage[T1]{fontenc}
  \usepackage[utf8]{inputenc}
  \usepackage{textcomp} % provide euro and other symbols
\else % if luatex or xetex
  \usepackage{unicode-math} % this also loads fontspec
  \defaultfontfeatures{Scale=MatchLowercase}
  \defaultfontfeatures[\rmfamily]{Ligatures=TeX,Scale=1}
\fi
\usepackage{lmodern}
\ifPDFTeX\else
  % xetex/luatex font selection
\fi
% Use upquote if available, for straight quotes in verbatim environments
\IfFileExists{upquote.sty}{\usepackage{upquote}}{}
\IfFileExists{microtype.sty}{% use microtype if available
  \usepackage[]{microtype}
  \UseMicrotypeSet[protrusion]{basicmath} % disable protrusion for tt fonts
}{}
\makeatletter
\@ifundefined{KOMAClassName}{% if non-KOMA class
  \IfFileExists{parskip.sty}{%
    \usepackage{parskip}
  }{% else
    \setlength{\parindent}{0pt}
    \setlength{\parskip}{6pt plus 2pt minus 1pt}}
}{% if KOMA class
  \KOMAoptions{parskip=half}}
\makeatother
% Make \paragraph and \subparagraph free-standing
\makeatletter
\ifx\paragraph\undefined\else
  \let\oldparagraph\paragraph
  \renewcommand{\paragraph}{
    \@ifstar
      \xxxParagraphStar
      \xxxParagraphNoStar
  }
  \newcommand{\xxxParagraphStar}[1]{\oldparagraph*{#1}\mbox{}}
  \newcommand{\xxxParagraphNoStar}[1]{\oldparagraph{#1}\mbox{}}
\fi
\ifx\subparagraph\undefined\else
  \let\oldsubparagraph\subparagraph
  \renewcommand{\subparagraph}{
    \@ifstar
      \xxxSubParagraphStar
      \xxxSubParagraphNoStar
  }
  \newcommand{\xxxSubParagraphStar}[1]{\oldsubparagraph*{#1}\mbox{}}
  \newcommand{\xxxSubParagraphNoStar}[1]{\oldsubparagraph{#1}\mbox{}}
\fi
\makeatother


\usepackage{longtable,booktabs,array}
\newcounter{none} % for unnumbered tables
\usepackage{calc} % for calculating minipage widths
% Correct order of tables after \paragraph or \subparagraph
\usepackage{etoolbox}
\makeatletter
\patchcmd\longtable{\par}{\if@noskipsec\mbox{}\fi\par}{}{}
\makeatother
% Allow footnotes in longtable head/foot
\IfFileExists{footnotehyper.sty}{\usepackage{footnotehyper}}{\usepackage{footnote}}
\makesavenoteenv{longtable}
\usepackage{graphicx}
\makeatletter
\newsavebox\pandoc@box
\newcommand*\pandocbounded[1]{% scales image to fit in text height/width
  \sbox\pandoc@box{#1}%
  \Gscale@div\@tempa{\textheight}{\dimexpr\ht\pandoc@box+\dp\pandoc@box\relax}%
  \Gscale@div\@tempb{\linewidth}{\wd\pandoc@box}%
  \ifdim\@tempb\p@<\@tempa\p@\let\@tempa\@tempb\fi% select the smaller of both
  \ifdim\@tempa\p@<\p@\scalebox{\@tempa}{\usebox\pandoc@box}%
  \else\usebox{\pandoc@box}%
  \fi%
}
% Set default figure placement to htbp
\def\fps@figure{htbp}
\makeatother


% definitions for citeproc citations
\NewDocumentCommand\citeproctext{}{}
\NewDocumentCommand\citeproc{mm}{%
  \begingroup\def\citeproctext{#2}\cite{#1}\endgroup}
\makeatletter
 % allow citations to break across lines
 \let\@cite@ofmt\@firstofone
 % avoid brackets around text for \cite:
 \def\@biblabel#1{}
 \def\@cite#1#2{{#1\if@tempswa , #2\fi}}
\makeatother
\newlength{\cslhangindent}
\setlength{\cslhangindent}{1.5em}
\newlength{\csllabelwidth}
\setlength{\csllabelwidth}{3em}
\newenvironment{CSLReferences}[2] % #1 hanging-indent, #2 entry-spacing
 {\begin{list}{}{%
  \setlength{\itemindent}{0pt}
  \setlength{\leftmargin}{0pt}
  \setlength{\parsep}{0pt}
  % turn on hanging indent if param 1 is 1
  \ifodd #1
   \setlength{\leftmargin}{\cslhangindent}
   \setlength{\itemindent}{-1\cslhangindent}
  \fi
  % set entry spacing
  \setlength{\itemsep}{#2\baselineskip}}}
 {\end{list}}
\usepackage{calc}
\newcommand{\CSLBlock}[1]{\hfill\break\parbox[t]{\linewidth}{\strut\ignorespaces#1\strut}}
\newcommand{\CSLLeftMargin}[1]{\parbox[t]{\csllabelwidth}{\strut#1\strut}}
\newcommand{\CSLRightInline}[1]{\parbox[t]{\linewidth - \csllabelwidth}{\strut#1\strut}}
\newcommand{\CSLIndent}[1]{\hspace{\cslhangindent}#1}

\ifLuaTeX
\usepackage[bidi=basic,shorthands=off]{babel}
\else
\usepackage[bidi=default,shorthands=off]{babel}
\fi
\ifLuaTeX
  \usepackage{selnolig} % disable illegal ligatures
\fi


\setlength{\emergencystretch}{3em} % prevent overfull lines

\providecommand{\tightlist}{%
  \setlength{\itemsep}{0pt}\setlength{\parskip}{0pt}}



 


% Get Math Fonts
\usepackage{fontspec}
\usepackage{unicode-math}
\setmathfont{Latin Modern Math}
\setmathfont[range={\mathscr,\mathbfscr,\mathbb}]{XITSMath-Regular}
[    Extension = .otf,
      BoldFont = XITSMath-Bold
]
\makeatletter
\@ifpackageloaded{caption}{}{\usepackage{caption}}
\AtBeginDocument{%
\ifdefined\contentsname
  \renewcommand*\contentsname{Table of contents}
\else
  \newcommand\contentsname{Table of contents}
\fi
\ifdefined\listfigurename
  \renewcommand*\listfigurename{List of Figures}
\else
  \newcommand\listfigurename{List of Figures}
\fi
\ifdefined\listtablename
  \renewcommand*\listtablename{List of Tables}
\else
  \newcommand\listtablename{List of Tables}
\fi
\ifdefined\figurename
  \renewcommand*\figurename{Figure}
\else
  \newcommand\figurename{Figure}
\fi
\ifdefined\tablename
  \renewcommand*\tablename{Table}
\else
  \newcommand\tablename{Table}
\fi
}
\@ifpackageloaded{float}{}{\usepackage{float}}
\floatstyle{ruled}
\@ifundefined{c@chapter}{\newfloat{codelisting}{h}{lop}}{\newfloat{codelisting}{h}{lop}[chapter]}
\floatname{codelisting}{Listing}
\newcommand*\listoflistings{\listof{codelisting}{List of Listings}}
\usepackage{amsthm}
\theoremstyle{definition}
\newtheorem{definition}{Definition}[section]
\theoremstyle{plain}
\newtheorem{theorem}{Theorem}[section]
\theoremstyle{remark}
\AtBeginDocument{\renewcommand*{\proofname}{Proof}}
\newtheorem*{remark}{Remark}
\newtheorem*{solution}{Solution}
\newtheorem{refremark}{Remark}[section]
\newtheorem{refsolution}{Solution}[section]
\makeatother
\makeatletter
\makeatother
\makeatletter
\@ifpackageloaded{caption}{}{\usepackage{caption}}
\@ifpackageloaded{subcaption}{}{\usepackage{subcaption}}
\makeatother

\newcommand{\Folner}[1][\vphantom{}]{{\Phi}_{#1}}
\newcommand{\FurstenbergFolner}[1][{\vphantom{}}]{{\Psi}_{#1}}
\newcommand{\Group}{{\Gamma}}
\newcommand{\GroupElement}{{\gamma}}
\newcommand{\Measure}{{\mu}}
\newcommand{\ShiftSpace}{{X}}

\usepackage{bookmark}
\IfFileExists{xurl.sty}{\usepackage{xurl}}{} % add URL line breaks if available
\urlstyle{same}
% fallback for those not using the hyperref driver hyperxmp:
\makeatletter
\@ifundefined{xmpquote}{\newcommand{\xmpquote}[1]{#1}}{}
\makeatother
\hypersetup{
  pdftitle={Erdős Cubes},
  pdflang={en-GB},
  colorlinks=true,
  linkcolor={blue},
  filecolor={Maroon},
  citecolor={Blue},
  urlcolor={Blue},
  pdfcreator={LaTeX via pandoc}}


\title{Erdős Cubes}
\author{}
\date{2026-09-26}
\begin{document}
\maketitle


Unless specified otherwise, let \((\ShiftSpace,\Measure,\Group)\) be the
Furstenberg system, \(\FurstenbergFolner\) a Følner sequence and \(a\)
be a point in \(\ShiftSpace\) defined by the countably-infinite discrete
amenable group \(\Group\), the Følner sequence \(\Folner\), and the
subset \(A\subset\Group\) where
\(A=\{\GroupElement\in\Group:T^{\GroupElement}a\in E\}\). Additionally,
let \([[k]]=\{0,1\}^k\) for some \(k\in\mathbb{N}\) and define
\(\ShiftSpace^{[[k]]}\) to be \(\prod_{i=1}^{2^k}\ShiftSpace\) and
\(T^{[[k]]}\) to be the corresponding \(2^k\)-fold transformation on
\(\ShiftSpace^{[[k]]}\).

We looked at ``Infinite Sumsets in Sets with Positive Density'\,'
(\citeproc{ref-kra2022infinite}{Kra et al., 2022b}) in our goal of
understanding Erdős cubes for the purpose of generalising the paper by
Host (\citeproc{ref-host2019short}{2019}) using more recent methods.

\begin{definition}[{Kra et al., 2022b, Definition
1.2}]\protect\hypertarget{def-2DErdosCube}{}\label{def-2DErdosCube}

We call \((x_{00},x_{01},x_{10},x_{11})\in\ShiftSpace^{[[2]]}\) a
\emph{2-dimensional Erdős cube} if there exists sequences
\(n\mapsto b(n),c(n)\) of distinct elements in \(\Group\) such that

\begin{align*}
(T\times T)^{c_1(n)}(x_{00},x_{01})&\rightarrow(x_{10},x_{11}),
\\(T\times T)^{c_2(n)}(x_{00},x_{10})&\rightarrow(x_{01},x_{11}).
\end{align*}

\end{definition}

Whilst understanding Erdős cubes, we also briefly compared it to Erdős
progressions to understand how these similar concepts differ.

\begin{definition}[{Kra et al., 2022a, Definition
2.1}]\protect\hypertarget{def-3ErdosProgression}{}\label{def-3ErdosProgression}

We call \((x_0,x_1,x_2)\in\ShiftSpace^3\) a \emph{3-term Erdős
progression} if there exists a sequence \(n\mapsto c(n)\) of distinct
elements in \(\Group\) such that
\[(T\times T)^{c(n)}(x_0,x_1)\rightarrow(x_1,x_2).\]

\end{definition}

We then worked through a basic example of showing that there exists an
Erdős cube in a Furstenberg dynamical system if and only if there exist
infinite sumsets within a positively dense set in a group.

\begin{theorem}[]\protect\hypertarget{thm-BCExample}{}\label{thm-BCExample}

The following statements are equivalent:

\begin{enumerate}
\def\labelenumi{\arabic{enumi}.}
\tightlist
\item
  There exists infinite sets \(B_1,B_2\subset\Group\) where
  \(B_1\cdot B_2\subset A\).
\item
  There exists an Erdős cube
  \((a,x_{01},x_{10},x_{11})\in\ShiftSpace^{[[2]]}\) where
  \(x_{11}\in E\).
\end{enumerate}

\end{theorem}

\begin{proof}
\leavevmode

\begin{enumerate}
\def\labelenumi{(\arabic{enumi})}
\tightlist
\item
  \(\Rightarrow\) (2): Assume we have infinite sets
  \(B_1,B_2\subset\Group\) where
  \[B_1\cdot B_2\subset A=\{\GroupElement\in\Group:T^ga\in E\}.\]
\end{enumerate}

Let \(n\mapsto b_1(n)\) and \(m\mapsto b_2(m)\) be infinite sequences of
distinct elements in \(B_1\) and \(B_2\), respectively. As
\(b_1(n)\cdot b_2(m)\in A\) for all \(n,m\in\mathbb{N}\), we have
\[T^{b_1(n)\cdot b_2(m)}a\in E \] for all \(n,m\in \mathbb{N}\). As
\(\ShiftSpace\) is compact, we use the property that every sequence has
a convergent subsequence to define subsequences
\((c_1(n))_{n\in\mathbb{N}}\) and \((c_2(m))_{m\in\mathbb{N}}\) such
that \begin{align*}
    \lim_{n\rightarrow\infty}T^{c_1(n)}a,
    \\\lim_{m\rightarrow\infty}T^{c_2(m)}a,
\end{align*} and \[\lim_{m\rightarrow\infty}T^{c_1(n)\cdot c_2(m)}a\]
exist and label these limits as \(x_{10}\), \(x_{01}\), and \(x_{11}\),
respectively. As \(T^{c_1(n)\cdot c_2(m)}a\in E\) for all
\(n,m\in\mathbb{N}\), we have \(x_{11}\in E\). Thus, we have shown the
existence of an Erdős cube where \(x_{11}\in E\).

\begin{enumerate}
\def\labelenumi{(\arabic{enumi})}
\setcounter{enumi}{1}
\tightlist
\item
  \(\Rightarrow\) (1): Assume we have an Erdős cube
  \((a,x_{01},x_{10},x_{11})\in\ShiftSpace^{[[2]]}\) where
  \(x_{11}\in E\). We construct subsequences \(n\mapsto b_1(n)\) and
  \(m\mapsto b_2(m)\) inductively. We start by choosing \(b_1(1)\) from
  \(\{c_1(n):n\in\mathbb{N}\}\) such that
  \[\lim_{m\rightarrow\infty}T^{b_1(1)\cdot c_2(m)}a\in E. \] Then, we
  choose \(b_2(1)\) from \(\{c_2(m):m\in\mathbb{N}\}\) such that
  \[\lim_{n\rightarrow\infty}T^{c_1(n)\cdot b_2(1)}a\in E. \] Next, we
  choose \(b_1(2)\neq b_1(1)\) from \(\{c_1(n):n\in\mathbb{N}\}\) such
  that \(b_1(2)\cdot b_2(1)\in\{\GroupElement\in\Group:T^ga\in E\}=A\)
  and \[\lim_{m\rightarrow\infty}T^{b_1(2)\cdot c_2(m)}a\in E. \] We
  then choose \(b_2(2)\neq b_2(1)\) from \(\{c_2(m):m\in\mathbb{N}\}\)
  such that \(b_1(1)\cdot b_2(2),\ b_1(2)\cdot b_2(2)\in A\) and
  \[\lim_{n\rightarrow\infty}T^{c_1(n)\cdot b_2(2)}a\in E. \] We repeat
  this alternating inductive construction by choosing
  \[b_1(i)\in\{c_1(n):n\in\mathbb{N}\}\setminus\{b_1(1),...,b_1(i-1)\}\]
  such that \(b_1(i)\cdot \{b_2(m):m<i\}\subset A\) and
  \[\lim_{m\rightarrow\infty}T^{b_1(i)\cdot c_2(m)}a\in E, \] then
  choosing
  \[b_2(i)\in\{c_2(m):m\in\mathbb{N}\}\setminus\{b_2(1),...,b_2(i-1)\}\]
  such that \(\{b_1(n):n\leq i\}\cdot b_2(i)\subset A\) and
  \[\lim_{n\rightarrow\infty}T^{c_1(n)\cdot b_2(i)}a\in E. \] After
  infinitely many steps, we are left with sequences
  \((b_1(n))_{n\in\mathbb{N}}\) and \((b_2(m))_{m\in\mathbb{N}}\) of
  distinct elements in \(\Group\) such that
  \(T^{b_1(n)\cdot b_2(m)}a\in E\) for all \(n,m\in\mathbb{N}\). As
  \(A=\{\GroupElement\in\Group:T^ga\in E\}\), we can conclude that
  \(b_1(n)\cdot b_2(m)\in A\) for all \(n,m\in\mathbb{N}\). By setting
  \(B_1=\{b_1(n):n\in\mathbb{N}\}\) and
  \(B_2=\{b_2(m):m\in\mathbb{N}\}\), we proven the existence of infinite
  subsets \(B_1,B_2\) such that \(B_1\cdot B_2\subset A\).
\end{enumerate}

\end{proof}

\protect\phantomsection\label{refs}
\begin{CSLReferences}{1}{0}
\bibitem[\citeproctext]{ref-host2019short}
Host, B. (2019). \emph{A short proof of a conjecture of erdös proved by
moreira, richter and robertson}. Available at:
\url{https://arxiv.org/abs/1904.09952}

\bibitem[\citeproctext]{ref-kra2022proof}
Kra, B. et al. (2022a). \emph{A proof of erdős's \(B+B+t\) conjecture}.
Available at: \url{https://arxiv.org/abs/2206.12377}

\bibitem[\citeproctext]{ref-kra2022infinite}
Kra, B. et al. (2022b). \emph{Infinite sumsets in sets with positive
density}. Available at: \url{https://arxiv.org/abs/2206.01786}

\end{CSLReferences}




\end{document}
